% =====================================================================
% 6. RESULTS
% =====================================================================
\section{Results}\label{sec:prdl-results}

\subsection{Evaluation methodology}

All numbers below are out-of-fold estimates from the nested cross-validation of
Section~\ref{sec:prdl-setup}, with 95\% bootstrap intervals in brackets. The coefficients, trees and
importances that are interpreted come from the final models fitted on all patients with the same
tuning procedure; their in-sample fit is not used as evidence of performance.

\subsection{Discussion}

\textbf{Regression: the profile explains part of the variance of age, and every model is
equivalent.} Predicting the mean age gives a mean absolute error of \RegBaseMAE{} years. All models
reduce it to between \RegMinMAE{} and \RegMaxMAE{} years and explain between \RegMinRTwo{} and
\RegMaxRTwo{} of the variance (Table~\ref{tab:prdl-reg}). The paired bootstrap confirms that every
model improves on the baseline. The \RegBest{} has the best point estimates (MAE \RegBestMAE{}
\RegBestMAECI{} years, $R^2$ \RegBestRTwo{} \RegBestRTwoCI), but its interval overlaps those of the
other \RegNTies{} models, and even the paired difference with linear regression includes zero
(\RegDiffOLSBest{} \RegDiffOLSBestCI{} years). By the rule fixed in advance, no model is declared
better: the additional flexibility of ensembles and kernels buys nothing that a linear model cannot
capture, which is expected when the predictors are binary and their effects nearly additive.

\begin{table}[H]
\centering
\caption{Regression of age at diagnosis: out-of-fold metrics with 95\% bootstrap intervals, and
$R^2$ per outer fold (mean $\pm$ SD).}
\label{tab:prdl-reg}
\resizebox{\textwidth}{!}{\input{tables/reg_model_comparison.tex}}
\end{table}

The interpretation is more interesting than the ranking. The multiple linear regression attributes
the largest effects to IDH1 (\OLSIDHOne{} years, 95\% CI \OLSIDHOneCI), IDH2 (\OLSIDHTwo{} years,
\OLSIDHTwoCI) and ATRX (\OLSATRX{} years); only \OLSNSigBH{} coefficients remain significant after
Benjamini--Hochberg correction, the maximum VIF is \OLSMaxVIF{}, and no patient has a Cook's distance
close to 1. The residual tests reject exact normality and homoscedasticity (Shapiro $p = \OLSShapirop$,
Breusch--Pagan $p = \OLSBPp$), so the coefficient intervals are approximate. One coefficient shows a
limitation of the data: the only American Indian or Alaska Native patient has leverage 1, so the model
fits that patient exactly and the coefficient (\OLSNative{} years, \OLSNativeCI) is meaningless.
Lasso, with the penalty chosen by cross-validation, sets \LassoNZero{} of \LassoNCoef{} coefficients
to zero and keeps the same leading predictors.

The out-of-fold predictions explain why the models agree (Figure~\ref{fig:prdl-oof}): they cluster
around two values. Their correlation with the grade is \RegCorrGrade, but with the IDH status (IDH1 or
IDH2 mutated) it is stronger, \RegCorrIDH{} (IDH-mutant patients are predicted younger), and Table~\ref{tab:prdl-reg-idh} separates the two: within GBM, the
\RegNGBMIDHmut{} IDH-mutant patients are predicted at \RegPredGBMIDHmut{} years on average, against
\RegPredGBMIDHwt{} for IDH-wildtype GBM, and within LGG the split is \RegPredLGGIDHmut{} vs
\RegPredLGGIDHwt{} years. The models predict the age from the IDH status, not from the grade; the
predictions correlate with the grade only because the grade is associated with IDH. A first draft of
this report read the two groups as the grade (Appendix~A); the causal analysis below quantifies how
much of the IDH--age association actually goes through the grade.

\begin{table}[H]
\centering
\caption{Mean observed and out-of-fold predicted age (years) of the \RegBest{} by grade and IDH status
(neither was an input of the model).}
\label{tab:prdl-reg-idh}
\small
\input{tables/reg_pred_by_grade_idh.tex}
\end{table}

\begin{figure}[H]
\centering
\includegraphics[width=0.9\textwidth]{reg_oof_predictions.png}
\caption{Out-of-fold predictions of the \RegBest{}: observed vs predicted age and residuals, coloured by
grade (not an input of the model).}
\label{fig:prdl-oof}
\end{figure}

\textbf{Classification: all models reach similar AUCs; simple models are as good as
ensembles.} Against a baseline with AUC \ClfAUCBase{} and accuracy \ClfBaseAcc, the models reach
AUCs between \ClfMinAUC{} and \ClfMaxAUC{} (Table~\ref{tab:prdl-clf}, Figure~\ref{fig:prdl-roc}). The
\ClfBest{} has the highest AUC (\ClfBestAUC{} \ClfBestAUCCI), with sensitivity \ClfBestSens{} and
specificity \ClfBestSpec{} at the default threshold, but the intervals of all \ClfNTies{} other models
overlap it, so by the pre-specified rule there is no winner. The paired comparison is more
informative: logistic regression (\ClfDiffLR) and XGBoost (\ClfDiffXGB) are indistinguishable from the
random forest, while naive Bayes (\ClfDiffNB), the SVM (\ClfDiffSVM) and the pruned tree
(\ClfDiffTree) are slightly but consistently worse. The two rules answer different questions, and
both are reported: overlapping marginal intervals say that the differences are small relative to the
uncertainty of each estimate; the paired test says that some of these small differences are
systematic. Either way, H1 is supported: a penalised logistic regression on these 23 predictors is as
good as any ensemble.

\begin{table}[H]
\centering
\caption{Classification of LGG vs GBM: out-of-fold metrics with 95\% bootstrap intervals (threshold
0.5 for sensitivity, specificity and balanced accuracy).}
\label{tab:prdl-clf}
\resizebox{\textwidth}{!}{\input{tables/clf_model_comparison.tex}}
\end{table}

\begin{figure}[H]
\centering
\includegraphics[width=0.55\textwidth]{clf_roc_curves.png}
\caption{Out-of-fold ROC curves of the six classifiers.}
\label{fig:prdl-roc}
\end{figure}

Three results deserve a critical reading.
\emph{(i) What the models learn.} Permutation importance ranks IDH1 far above every other predictor
(mean AUC drop \ClfPermIDH{}), followed by age; the pruned tree (depth \ClfTreeDepth, \ClfTreeLeaves{}
leaves) first splits on IDH1 and then on IDH2, TP53 and age. The models reproduce the IDH axis found
without labels in the EDA, which is consistent with the role of IDH in the WHO classification. The
\ClfFN{} GBM patients predicted as LGG and the \ClfFP{} LGG patients predicted as GBM by the
\ClfBest{} cannot be explained from these data alone; part of them may be label noise from the
pre-2021 diagnoses (Section 2.1).
\emph{(ii) Separation and the meaning of coefficients.} The final logistic regression chose an
almost unpenalised L1 model ($C = \LRFinalC$) although the outer folds chose smaller values: the
inner-CV surface is flat (\LRFlatWithin{} of \LRFlatTotal{} configurations within 0.005 AUC of the
best), so the choice among them is close to arbitrary. With a weak penalty, NOTCH1, mutated only in
LGG patients (in TCGA, a mutation of 1p/19q-codeleted gliomas \cite{tcga2015lgg}), receives a coefficient (\LRCoefNOTCH{} log-odds per SD) comparable to that of IDH1
(\LRCoefIDH): a direct consequence of quasi-separation \cite{albert1984}, not evidence that NOTCH1 is
as important as IDH1. The odds ratio of age, \LRORAge{} per standard deviation, is more stable.
\emph{(iii) Ranking vs calibration.} Naive Bayes reaches a competitive AUC, but \NBExtreme\% of its
out-of-fold probabilities are below 0.02 or above 0.98, against \LRExtreme\% for logistic regression.
Because it counts the evidence of correlated mutations (IDH1, ATRX, TP53) several times, its
probabilities are over-confident even when its ranking is good \cite{hand2001}. If the probabilities
were to be used, for example to choose a clinical threshold, this difference would matter more than
the AUC.

\textbf{Comparison with previous work.} On the same 839 TCGA patients, \cite{tasci2022} reported an
accuracy of 87.606\% with a voting ensemble and feature selection, and \cite{sanchezmarques2024}
0.882 with oversampling. The accuracies obtained here (between \ClfMinAcc{} and \ClfMaxAcc, MLLB report) are of the same order with standard models and no resampling. A strict comparison is not
possible because the validation schemes differ. In \cite{sanchezmarques2024} the hyperparameters were
``fine-tuned [...] using an 80-20 stratified holdout setting'' and performance was then evaluated with
10-fold cross-validation; if the tuning split is drawn from the same patients, the tuned values have
seen the evaluation folds, and the estimate is optimistic \cite{varma2006,cawley2010}.
In \cite{tasci2022}, instead, no tuning was done: ``the default values were assigned as the
corresponding parameter values for the utilized classifiers'', so their estimate is not inflated by
hyperparameter selection, but their classifiers were not tuned either. The nested cross-validation
used here tunes every model while keeping the evaluation folds out of the tuning, at the cost of
slightly lower but honest numbers.

\textbf{Multiclass: the mutations recognise the families, except the mixed gliomas.} Against a baseline
macro AUC of \McBaseAUC, the four models reach between \McMinAUC{} and \McMaxAUC{}
(Table~\ref{tab:prdl-mc}). The \McBest{} has the highest value (\McBestAUC; balanced accuracy
\McBestBAcc, macro F1 \McBestFOne), the intervals of the \McNTies{} other models overlap it, and in the
paired comparison only the \McWorse{} is worse. The probability of glioblastoma ranks the grade almost as
well as the binary models (AUC \McGradeAUC), so the three extra classes cost nothing for the grade. The
confusion matrix (Figure~\ref{fig:prdl-mc}) shows where the information lies:
\begin{itemize}
  \item \textbf{Glioblastoma} is recognised best (recall \McRecallGBM, one-vs-rest AUC \McAUCGBM),
        followed by \textbf{oligodendroglioma} (\McRecallOligo, \McAUCOligo) and \textbf{astrocytoma}
        (\McRecallAstro, \McAUCAstro).
  \item \McAstroAsGBM{} astrocytomas are predicted as glioblastoma; \McAstroGBMwt{} of them are
        IDH-wildtype, and they are all the IDH-wildtype astrocytomas (\McAstroWtTotal) of the cohort.
        In the 2021 classification, for an IDH-wildtype diffuse astrocytic glioma the ``presence of 1 or
        more of 3 genetic parameters (TERT promoter mutation, EGFR gene amplification, combined gain of
        entire chromosome 7 and loss of entire chromosome 10 [+7/$-$10]) appears sufficient to assign''
        the diagnosis of glioblastoma, IDH-wildtype \cite{louis2021}. These markers are not in the data,
        so how many of the \McAstroGBMwt{} patients would be reclassified cannot be checked, but these
        ``errors'' are at least partly the outdated labels of Section~2.1.
  \item No patient is predicted as \textbf{mixed glioma} (recall \McRecallMixed): \McMixedAsAstro\% are
        predicted as astrocytoma, \McMixedAsOligo\% as oligodendroglioma and \McMixedAsGBM\% as
        glioblastoma, and their mean predicted probability is split between the two pure families
        (Figure~\ref{fig:prdl-mc}, right). Their mutation profile is that of one of the two pure
        families. This is what the 2016 WHO classification anticipated: ``the diagnosis of
        oligoastrocytoma is strongly discouraged. Nearly all tumors with histological features suggesting
        both an astrocytic and an oligodendroglial component can be classified as either astrocytoma or
        oligodendroglioma using genetic testing'' \cite{louis2016}. Here, point mutations alone already
        push most of them towards one of the two families.
\end{itemize}
The unsupervised analysis of Section~\ref{sec:prdl-eda} anticipated this structure without labels: the
CIC/FUBP1 cluster is \McCThreeOligo\% oligodendrogliomas but also contains \McCThreeMixed{} mixed
gliomas, and the TP53/ATRX cluster is \McCTwoAstro\% astrocytomas with \McCTwoMixed{} mixed gliomas.

\begin{table}[H]
\centering
\caption{Multiclass classification of the histological family: out-of-fold metrics with 95\% bootstrap
intervals (\McNAstro{} astrocytomas, \McNMixed{} mixed gliomas, \McNOligo{} oligodendrogliomas, \McNGBM{}
glioblastomas).}
\label{tab:prdl-mc}
\resizebox{\textwidth}{!}{\input{tables/mc_model_comparison.tex}}
\end{table}

\begin{figure}[H]
\centering
\includegraphics[width=\textwidth]{mc_confusion.png}
\caption{\McBest: out-of-fold confusion matrix (colour: share of each true family) and mean predicted
probability of each family by true family.}
\label{fig:prdl-mc}
\end{figure}

\textbf{What the winning models use: SHAP.} SHAP values were computed for the three binary models
that the paired comparison could not separate, and for the multiclass \ShapMcModel{}
(Figures~\ref{fig:prdl-shap} and~\ref{fig:prdl-shap-mc}). Their global ranking agrees with the
permutation importance (Spearman $\rho$ = \ShapRhoRF{} for the random forest, \ShapRhoLR{} for logistic
regression and \ShapRhoXGB{} for XGBoost), and IDH1 takes \ShapIDHRF\% of the total mean $|$SHAP$|$ of
the random forest and \ShapIDHXGB\% of XGBoost. SHAP adds three things that the permutation importance
could not show:
\begin{itemize}
  \item \textbf{The direction and size of each effect, patient by patient.} NOTCH1 is the second input of
        logistic regression (\ShapNOTCHLR\% of its total) but negligible in the random forest
        (\ShapNOTCHRF\%): the few NOTCH1-mutated patients receive a very large contribution towards LGG,
        the quasi-separation effect discussed above made visible.
  \item \textbf{The shape of the effect of age.} In logistic regression it is a straight line by
        construction; the tree ensembles learn a curve that is almost flat for young patients and rises
        steeply in middle age, changing sign at about \ShapAgeRF{} years in the random forest and
        \ShapAgeXGB{} in XGBoost.
  \item \textbf{What separates the families.} In the multiclass forest, CIC weighs most for
        oligodendroglioma (\ShapCICOligo\% of that class's total, against \ShapCICAstro\% for
        astrocytoma), and ATRX and age more for astrocytoma (\ShapATRXAstro\% and \ShapAgeAstro\%, against
        \ShapATRXOligo\% and \ShapAgeOligo\%), the molecular split described for the TCGA lower-grade
        gliomas \cite{tcga2015lgg}.
\end{itemize}
The ranking agrees with the previous explainability analysis of these data among the sources reviewed:
on the same patients, Sánchez-Marqués et al.\ combined four feature-ranking methods, permutation
importance and SHAP values, and report that the ``IDH1 biomarker was the most important attribute
contributing to the target variable (i.e., glioma grade), regardless of the model used. The second
most important feature was Age at diagnosis'' \cite{sanchezmarques2024}. The same order appears here in
the random forest and XGBoost; logistic regression is the exception, for the quasi-separation reason
above. Tasci et al.\ used feature selection but do not report a ranking of the features
\cite{tasci2022}.

\begin{figure}[H]
\centering
\includegraphics[width=\textwidth]{shap_binary_summary.png}
\caption{SHAP summary of the three tied binary models: one dot per patient, contribution towards GBM
(colour: value of the input, grey = 0 or young, violet = 1 or old). Units: probability for the random
forest, log-odds for logistic regression and XGBoost.}
\label{fig:prdl-shap}
\end{figure}

\begin{figure}[H]
\centering
\includegraphics[width=\textwidth]{shap_multiclass.png}
\caption{Multiclass \ShapMcModel: share of mean $|$SHAP$|$ of each input per family, and contributions
towards astrocytoma and oligodendroglioma.}
\label{fig:prdl-shap-mc}
\end{figure}

\subsection{Causal analysis}\label{sec:prdl-causal}

Figure~\ref{fig:prdl-dag} shows the diagram assumed for Q5. IDH-mutant patients are younger within both
grades (LGG: \CsAgeLGGmut{} vs \CsAgeLGGwt{} years; GBM: \CsAgeGBMmut{} vs \CsAgeGBMwt{} years, with
only \CsNGBMmut{} IDH-mutant GBM), which already suggests that the grade does not carry the whole
association. Adjusted for sex and race, the total effect of an IDH mutation on the age at diagnosis is
\CsTotal{} years (OLS with HC3 standard errors \cite{hayes2007}; bootstrap intervals). Adjusting also
for the grade leaves a direct effect of \CsDirect{} years; the indirect effect through the grade is
\CsIndirect{} years, \CsPropPct\% of the total (Table~\ref{tab:prdl-med}). The identity
$c = c' + ab$ holds exactly, so the result rests on the assumptions, not on the arithmetic.

\begin{figure}[H]
\centering
\begin{minipage}[t]{0.47\textwidth}\centering
\includegraphics[width=\textwidth]{causal_dag.png}
\caption{Assumed causal diagram (dashed: unmeasured).}
\label{fig:prdl-dag}
\end{minipage}\hfill
\begin{minipage}[t]{0.5\textwidth}\centering
\includegraphics[width=\textwidth]{causal_adjustment_sets.png}
\caption{Coefficient of IDH on age for different adjustment sets.}
\label{fig:prdl-adj}
\end{minipage}
\end{figure}

\begin{table}[H]
\centering
\caption{Mediation decomposition of the IDH--age association (adjusted for sex and race; 95\% bootstrap
intervals, 2000 resamples).}
\label{tab:prdl-med}
\small
\input{tables/causal_mediation.tex}
\end{table}

Three readings follow:
\begin{itemize}
  \item \textbf{It corrects H2.} Under the diagram, most of the association between IDH and age does
        not go through the binary grade, in agreement with the regression predictions above.
  \item \textbf{What is adjusted for changes the answer} (Figure~\ref{fig:prdl-adj}). Adjusting for the
        histological family instead of the grade gives \CsFamily{} years, and adding the 18 other
        mutations \CsGenes{} years. The family is itself a consequence of IDH, so adjusting for it is
        over-adjustment for the total effect; whether the other mutations are confounders or mediators
        the diagram does not settle. Adding variables is not ``more control'': the diagram decides
        \cite{huntingtonklein2021,pearl2009}. For the prediction models above the same question does
        not arise, because there any variable that carries information is useful.
  \item \textbf{The direct effect is the fragile part.} If an unmeasured variable affects both the grade
        and the age at diagnosis (U in Figure~\ref{fig:prdl-dag}, for example the copy-number markers
        missing from the data), adjusting for the grade opens the path IDH $\to$ grade $\leftarrow$ U
        $\to$ age, and $c'$ is biased; the total effect does not need that assumption. The decomposition
        also assumes that the effect of IDH is the same in both grades, but the IDH $\times$ grade
        interaction is \CsInter{} years, so it is probably larger in GBM, although that estimate rests
        on \CsNGBMmut{} patients.
\end{itemize}
